\section{Глава 2. Мультиагентная централизованная задача}

Следующий этап работы - решение задачи передвижения нескольких самолётов по проезжей части аэропорта,
от ворот до ближайших выходов.
В первую очередь будет опробован централизованный подход, где агентом выступает "диспетчер", 
организующий движение каждого самолёта на карте одновременно.
На вход агент также получает единовременно одно состояние, описывающие положения, направления и карты проезжей части 
с отмеченными пунктами назначения каждого самолёта, и на основе всей информации от каждого самолёта принимается решение об их перемещениях. 
По сравнению с обучающим процессом из предыдущей главы, задача усложняется большим количеством принимаемых решений и менее явным возвращаемым 
значением награды - в случае централизованного подхода, награда равна сумме наград каждого подагента \eqref{eq:r_cen}, то есть оценивается влияние
диспетчера на всю систему в целом. 
Вследствие этого агенту сложнее определить, какие действия на каком самолёте оказались более успешными.

\begin{equation}
\label{eq:r_cen}
R_t = \sum_i r_t^i
\end{equation}

Сценарий эксперимента с применением шока остаётся неизменным по сравнению с одно-агентной задачей.

\subsection{Архитектура сети}

\subsection{Результаты}

